% Physical pages: 97--104 (printed pages 74--81).
% Figures: none.
% Uncertainties: none.

\subsection{Space Inversion and Time-Reversal}

We saw in Section 2.3 that any homogeneous Lorentz transformation is either
proper and orthochronous (i.e., $\operatorname{Det}\Lambda=+1$ and
$\Lambda^0{}_0\geq+1$) or else equal to a proper orthochronous transformation
times either $\mathcal{P}$ or $\mathcal{T}$ or $\mathcal{P}\mathcal{T}$, where
$\mathcal{P}$ and $\mathcal{T}$ are the space inversion and time-reversal
transformations
\[
\mathcal{P}^{\mu}{}_{\nu}=
\begin{pmatrix}
1&0&0&0\\
0&-1&0&0\\
0&0&-1&0\\
0&0&0&-1
\end{pmatrix},
\qquad
\mathcal{T}^{\mu}{}_{\nu}=
\begin{pmatrix}
-1&0&0&0\\
0&1&0&0\\
0&0&1&0\\
0&0&0&1
\end{pmatrix}.
\]
It used to be thought self-evident that the fundamental multiplication rule
of the Poincar\'e group
\[
U(\bar\Lambda,\bar a)U(\Lambda,a)
=U(\bar\Lambda\Lambda,\bar\Lambda a+\bar a)
\]
would be valid even if $\Lambda$ and/or $\bar\Lambda$ involved factors of
$\mathcal{P}$ or $\mathcal{T}$ or $\mathcal{P}\mathcal{T}$. In particular, it
was believed that there are operators corresponding to $\mathcal{P}$ and
$\mathcal{T}$ themselves:
\[
P\equiv U(\mathcal{P},0),
\qquad
T\equiv U(\mathcal{T},0)
\]
such that
\begin{align}
P U(\Lambda,a)P^{-1}
  &=U(\mathcal{P}\Lambda\mathcal{P}^{-1},\mathcal{P}a),
  \tag{2.6.1}\label{eq:2.6.1}\\
T U(\Lambda,a)T^{-1}
  &=U(\mathcal{T}\Lambda\mathcal{T}^{-1},\mathcal{T}a)
  \tag{2.6.2}\label{eq:2.6.2}
\end{align}
for any proper orthochronous Lorentz transformation $\Lambda^\mu{}_{\nu}$,
and translation $a^\mu$. These transformation rules incorporate most of what
is meant when we say that $P$ or $T$ are `conserved'.

In 1956--57 it became understood~\hyperref[ref:2.8]{[8]} that this is true for $P$
only in the approximation in which one ignores the effects of weak
interactions, such as those that produce nuclear beta decay. Time-reversal
survived for a while, but in 1964 there appeared indirect
evidence~\hyperref[ref:2.9]{[9]} that these properties of $T$ are also only
approximately satisfied. (See Section 3.3.) In what follows, we will make
believe that operators $P$ and $T$ satisfying Eqs.~\eqref{eq:2.6.1} and \eqref{eq:2.6.2}
actually exist, but it should be kept in mind that this is only an
approximation.

Let us apply Eqs.~\eqref{eq:2.6.1} and \eqref{eq:2.6.2} in the case of an infinitesimal
transformation, i.e.,
\[
\Lambda^\mu{}_{\nu}=\delta^\mu{}_{\nu}+\omega^\mu{}_{\nu},
\qquad
a^\mu=\epsilon^\mu
\]
with $\omega_{\mu\nu}=-\omega_{\nu\mu}$ and $\epsilon_\mu$ both infinitesimal.
Using \eqref{eq:2.4.3}, and equating coefficients of $\omega_{\rho\sigma}$ and
$\epsilon_\rho$ in Eqs.~\eqref{eq:2.6.1} and \eqref{eq:2.6.2}, we obtain the $P$ and $T$
transformation properties of the Poincar\'e generators
\begin{align}
P iJ^{\rho\sigma}P^{-1}
  &=i\mathcal{P}_\mu{}^\rho\mathcal{P}_\nu{}^\sigma J^{\mu\nu},
  \tag{2.6.3}\label{eq:2.6.3}\\
P iP^\rho P^{-1}
  &=i\mathcal{P}_\mu{}^\rho P^\mu,
  \tag{2.6.4}\label{eq:2.6.4}\\
T iJ^{\rho\sigma}T^{-1}
  &=i\mathcal{T}_\mu{}^\rho\mathcal{T}_\nu{}^\sigma J^{\mu\nu},
  \tag{2.6.5}\label{eq:2.6.5}\\
T iP^\rho T^{-1}
  &=i\mathcal{T}_\mu{}^\rho P^\mu.
  \tag{2.6.6}\label{eq:2.6.6}
\end{align}
This is much like Eqs.~\eqref{eq:2.4.8} and \eqref{eq:2.4.9}, except that we have not cancelled
factors of $i$ on both sides of these equations, because at this point we have
not yet decided whether $P$ and $T$ are linear and unitary or antilinear and
antiunitary.

The decision is an easy one. Setting $\rho=0$ in Eq.~\eqref{eq:2.6.4} gives
\[
P iH P^{-1}=iH,
\]
where $H\equiv P^0$ is the energy operator. If $P$ were antiunitary and
antilinear then it would anticommute with $i$, so $PHP^{-1}=-H$. But then for
any state $\ket{\Psi}$ of energy $E>0$, there would have to be another state
$P^{-1}\ket{\Psi}$ of energy $-E<0$. There are no states of negative energy
(energy less than that of the vacuum), so we are forced to choose the other
alternative: $P$ \emph{is linear and unitary, and commutes rather than
anticommutes with $H$}.

On the other hand, setting $\rho=0$ in Eq.~\eqref{eq:2.6.6} yields
\[
T iH T^{-1}=-iH.
\]
If we supposed that $T$ is linear and unitary then we could simply cancel the
$i$s, and find $THT^{-1}=-H$, with the again disastrous conclusion that for
any state $\ket{\Psi}$ of energy $E$ there is another state
$T^{-1}\ket{\Psi}$ of energy $-E$. To avoid this, we are forced here to
conclude that $T$ \emph{is antilinear and antiunitary}.

\begin{sloppypar}
Now that we have decided that $P$ is linear and $T$ is antilinear, we can
conveniently rewrite Eqs.~\eqref{eq:2.6.3}--\eqref{eq:2.6.6} in terms of the generators
\eqref{eq:2.4.15}--\eqref{eq:2.4.17} in a three-dimensional notation
\end{sloppypar}
\begin{align}
P\mathbf{J}P^{-1}&=+\mathbf{J}, & \tag{2.6.7}\label{eq:2.6.7}\\
P\mathbf{K}P^{-1}&=-\mathbf{K}, & \tag{2.6.8}\label{eq:2.6.8}\\
P\mathbf{P}P^{-1}&=-\mathbf{P}, & \tag{2.6.9}\label{eq:2.6.9}\\
T\mathbf{J}T^{-1}&=-\mathbf{J}, & \tag{2.6.10}\label{eq:2.6.10}\\
T\mathbf{K}T^{-1}&=+\mathbf{K}, & \tag{2.6.11}\label{eq:2.6.11}\\
T\mathbf{P}T^{-1}&=-\mathbf{P}, & \tag{2.6.12}\label{eq:2.6.12}
\end{align}
and, as shown before,
\begin{equation}
PHP^{-1}=THT^{-1}=H.
\tag{2.6.13}\label{eq:2.6.13}
\end{equation}
It is physically sensible that $P$ should preserve the sign of $\mathbf{J}$,
because at least the orbital part is a vector product
$\mathbf{r}\cdot\mathbf{p}$ of two vectors, both of which change sign under
an inversion of the spatial coordinate system. On the other hand, $T$
reverses $\mathbf{J}$, because after time-reversal an observer will see all
bodies spinning in the opposite direction. Note by the way that Eq.~\eqref{eq:2.6.10}
is consistent with the angular-momentum commutation relations
$\mathbf{J}\cdot\mathbf{J}=i\mathbf{J}$, because $T$ reverses not only
$\mathbf{J}$, but also $i$. The reader can easily check that Eqs.
\eqref{eq:2.6.7}--\eqref{eq:2.6.13} are consistent with all the commutation relations
\eqref{eq:2.4.18}--\eqref{eq:2.4.24}.

Let us now consider what $P$ and $T$ do to one-particle states:

\medskip
\noindent\textbf{$P:M>0$}

The one-particle states $\ket{k,\sigma}$ are defined as eigenvectors of
$\mathbf{P}$, $H$, and $J_3$ with eigenvalues $0$, $M$, and $\sigma$,
respectively. From Eqs.~\eqref{eq:2.6.7}, \eqref{eq:2.6.9}, and \eqref{eq:2.6.13}, we see that the same
must be true of the state $P\ket{k,\sigma}$, and therefore (barring
degeneracies) these states can only differ by a phase
\[
P\ket{k,\sigma}=\eta_\sigma\ket{k,\sigma}
\]
with a phase factor ($|\eta_\sigma|=1$) that may or may not depend on the spin
$\sigma$. To see that $\eta_\sigma$ is $\sigma$-independent, we note from
\eqref{eq:2.5.8}, \eqref{eq:2.5.20}, and \eqref{eq:2.5.21} that
\begin{equation}
(J_1\mathbin{\pm}iJ_2)\ket{k,\sigma}
=\sqrt{(j\mp\sigma)(j\pm\sigma+1)}\ket{k,\sigma\mathbin{\pm}1},
\tag{2.6.14}\label{eq:2.6.14}
\end{equation}
where $j$ is the particle's spin. Operating on both sides with $P$, we find
\[
\eta_\sigma=\eta_{\sigma\pm1}
\]
and so $\eta_\sigma$ is actually independent of $\sigma$. We therefore write
\begin{equation}
P\ket{k,\sigma}=\eta\ket{k,\sigma}
\tag{2.6.15}\label{eq:2.6.15}
\end{equation}
with $\eta$ a phase, known as the \emph{intrinsic parity}, that depends only
on the species of particle on which $P$ acts.

To get to finite momentum states, we must apply the unitary operator
$U(L(p))$ corresponding to the `boost' \eqref{eq:2.5.24}:
\[
\ket{p,\sigma}=\sqrt{M/p^0}\,U(L(p))\ket{k,\sigma}.
\]
We note that
\[
\mathcal{P}L(p)\mathcal{P}^{-1}=L(\mathcal{P}p),
\qquad
\mathcal{P}p=(p^0,-\mathbf{p}),
\]
so using Eqs.~\eqref{eq:2.6.1} and \eqref{eq:2.6.15}, we have
\[
P\ket{p,\sigma}
=\sqrt{M/p^0}\,U(L(\mathcal{P}p))\eta\ket{k,\sigma}
\]
or in other words
\begin{equation}
P\ket{p,\sigma}=\eta\ket{\mathcal{P}p,\sigma}.
\tag{2.6.16}\label{eq:2.6.16}
\end{equation}

\medskip
\noindent\textbf{$T:M>0$}

From Eqs.~\eqref{eq:2.6.10}, \eqref{eq:2.6.12}, and \eqref{eq:2.6.13}, we see that the effect of $T$ on
the zero-momentum one-particle state $\ket{k,\sigma}$ is to yield a state with
\[
\begin{aligned}
\mathbf{P}(T\ket{k,\sigma})&=0,\\
H(T\ket{k,\sigma})&=M(T\ket{k,\sigma}),\\
J_3(T\ket{k,\sigma})&=-\sigma(T\ket{k,\sigma}),
\end{aligned}
\]
and so
\[
T\ket{k,\sigma}=\zeta_\sigma\ket{k,-\sigma},
\]
where $\zeta_\sigma$ is a phase factor. Applying the operator $T$ to
\eqref{eq:2.6.14}, and recalling that $T$ anticommutes with $\mathbf{J}$ \emph{and}
$i$, we find
\[
(-J_1\mathbin{\pm}iJ_2)\zeta_\sigma\ket{k,-\sigma}
=\sqrt{(j\mp\sigma)(j\pm\sigma+1)}\,
 \zeta_{\sigma\pm1}\ket{k,-\sigma\mp1}.
\]
Using Eq.~\eqref{eq:2.6.14} again on the left, we see that the square-root factors
cancel, and so
\[
-\zeta_\sigma=\zeta_{\sigma\pm1}.
\]
We write the solution as $\zeta_\sigma=\zeta(-1)^{j-\sigma}$, with $\zeta$
another phase that depends only on the species of particle:
\begin{equation}
T\ket{k,\sigma}=\zeta(-1)^{j-\sigma}\ket{k,-\sigma}.
\tag{2.6.17}\label{eq:2.6.17}
\end{equation}
However, unlike the `intrinsic parity' $\eta$, the time-reversal phase $\zeta$
has no physical significance. This is because we can redefine the
one-particle states by a change of phase
\[
\ket{k,\sigma}\longrightarrow
\ket{k,\sigma}'=\zeta^{1/2}\ket{k,\sigma}
\]
in such a way that the phase $\zeta$ is eliminated from the transformation
rule
\[
T\ket{k,\sigma}'
=\zeta^{*1/2}T\ket{k,\sigma}
=\zeta^{*1/2}\zeta(-1)^{j-\sigma}\ket{k,-\sigma}
=(-1)^{j-\sigma}\ket{k,-\sigma}'.
\]
In what follows we will keep the arbitrary phase $\zeta$ in Eq.~\eqref{eq:2.6.17},
just to keep open our options in choosing the phase of the one-particle states,
but it should be kept in mind that this phase is of no real importance.

To deal with states of finite momentum, we again apply the `boost' \eqref{eq:2.5.24}.
Note that
\[
\mathcal{T}L(p)\mathcal{T}^{-1}=L(\mathcal{P}p),
\qquad
\mathcal{P}p=(p^0,-\mathbf{p}).
\]
(That is, changing the sign of each element of $L^\mu{}_{\nu}$ with an odd
number of \emph{time}-indices is the same as changing the signs of elements
with an odd number of \emph{space}-indices.) Using Eqs.~\eqref{eq:2.6.2} and \eqref{eq:2.5.5},
we have then
\begin{equation}
T\ket{p,\sigma}
=\zeta(-1)^{j-\sigma}\ket{\mathcal{P}p,-\sigma}.
\tag{2.6.18}\label{eq:2.6.18}
\end{equation}

\medskip
\noindent\textbf{$P:M=0$}

Acting on a state $\ket{k,\sigma}$ that is defined as an eigenvector of
$P^\mu$ with eigenvalue $k^\mu=(\kappa,0,0,\kappa)$ and an eigenvector of
$J_3$ with eigenvalue $\sigma$, the parity operator $P$ yields a state with
four-momentum $(\mathcal{P}k)^\mu=(\kappa,0,0,-\kappa)$ and $J_3$ equal to
$\sigma$. Thus it takes a state of helicity (the component of spin along the
direction of motion) $\sigma$ into one of helicity $-\sigma$. As mentioned
earlier, this shows that the existence of a space-inversion symmetry requires
that any species of massless particle with non-zero helicity must be
accompanied with another of opposite helicity. Because $P$ does not leave the
standard momentum invariant, it is convenient to consider instead the
operator $U(R_2^{-1})P$, where $R_2$ is a rotation that also takes $k$ to
$\mathcal{P}k$, conveniently chosen as a rotation by $-180^\circ$ around the
two-axis
\begin{equation}
U(R_2)=\exp(-i\pi J_2).
\tag{2.6.19}\label{eq:2.6.19}
\end{equation}
Since $U(R_2^{-1})$ reverses the sign of $J_3$, we have
\begin{equation}
U(R_2^{-1})P\ket{k,\sigma}=\eta_\sigma\ket{k,-\sigma}
\tag{2.6.20}\label{eq:2.6.20}
\end{equation}
with $\eta_\sigma$ a phase factor. Now, $R_2^{-1}\mathcal{P}$ commutes with the
Lorentz `boost' \eqref{eq:2.5.45}, and $\mathcal{P}$ commutes with the rotation
$R(\hat{\mathbf{p}})$ which takes the three-direction into the direction of
$\mathbf{p}$, so by operating on \eqref{eq:2.5.5} with $P$, we find for a general
four-momentum $p^\mu$
\[
\begin{aligned}
P\ket{p,\sigma}
&=\sqrt{\frac{\kappa}{p^0}}\,
U\left(R(\hat{\mathbf{p}})R_2B\left(\frac{|\mathbf{p}|}{\kappa}\right)\right)
U(R_2^{-1})P\ket{k,\sigma}\\
&=\sqrt{\frac{\kappa}{p^0}}\,\eta_\sigma
U\left(R(\hat{\mathbf{p}})R_2B\left(\frac{|\mathbf{p}|}{\kappa}\right)\right)
\ket{k,-\sigma}.
\end{aligned}
\]
Note that $R(\hat{\mathbf{p}})R_2$ is a rotation that takes the three-axis into
the direction of $-\hat{\mathbf{p}}$, but
$U(R(\hat{\mathbf{p}})R_2)$ is not quite equal to
$U(R(-\hat{\mathbf{p}}))$. According to \eqref{eq:2.5.47},
\[
U(R(-\hat{\mathbf{p}}))
=\exp\bigl(i(\phi\pm\pi)J_3\bigr)
 \exp\bigl(i(\pi-\theta)J_2\bigr)
\]
with azimuthal angle chosen as $\phi+\pi$ or $\phi-\pi$ according to whether
$0\leq\phi<\pi$ or $\pi\leq\phi<2\pi$, so that it remains in the range of
$0$ to $2\pi$. Then
\[
\begin{aligned}
U^{-1}(R(-\hat{\mathbf{p}}))U(R(\hat{\mathbf{p}})R_2)
&=\exp\bigl(-i(\pi-\theta)J_2\bigr)\\
&\quad\cdot\exp\bigl(-i(\phi\pm\pi)J_3\bigr)
\exp(i\phi J_3)\exp(i\theta J_2)\exp(-i\pi J_2)\\
&=\exp\bigl(-i(\pi-\theta)J_2\bigr)
\exp(\mp i\pi J_3)
\exp\bigl(-i(\pi-\theta)J_2\bigr).
\end{aligned}
\]
But a rotation of $\pm180^\circ$ around the three-axis reverses the sign of
$J_2$, so
\begin{equation}
U(R(\hat{\mathbf{p}})R_2)
=U(R(-\hat{\mathbf{p}}))\exp(\pm i\pi J_3).
\tag{2.6.21}\label{eq:2.6.21}
\end{equation}
Also, $R(-\hat{\mathbf{p}})B(|\mathbf{p}|/\kappa)$ is just the standard boost
$L(\mathcal{P}p)$ in the direction $\mathcal{P}p=(p^0,-\mathbf{p})$. We have
then finally
\begin{equation}
P\ket{p,\sigma}
=\eta_\sigma\exp(\mp i\pi\sigma)\ket{\mathcal{P}p,-\sigma}
\tag{2.6.22}\label{eq:2.6.22}
\end{equation}
with the phase $-\pi\sigma$ or $+\pi\sigma$ according to whether the
two-component of $\mathbf{p}$ is positive or negative, respectively. This
peculiar change of sign in the operation of parity for massless particles of
half-integer spin is due to the convention adopted in Eq.~\eqref{eq:2.5.47} for the
rotation used to define massless particle states of arbitrary momentum.
Because the rotation group is not simply connected, some discontinuity of
this sort is unavoidable.

\medskip
\noindent\textbf{$T:M=0$}

Acting on the state $\ket{k,\sigma}$, which has values
$k^\mu=(\kappa,0,0,\kappa)$ and $\sigma$ for $P^\mu$ and $J_3$, the
time-reversal operator $T$ yields a state which has values
$(\mathcal{P}k)^\mu=(\kappa,0,0,-\kappa)$ and $-\sigma$ for $P^\mu$ and $J_3$.
Thus $T$ does not change the helicity $\mathbf{J}\cdot\hat{\mathbf{k}}$, and
by itself has nothing to say about whether massless particles of one helicity
$\sigma$ are accompanied with others of helicity $-\sigma$. Because $T$ like
$P$ does not leave the standard four-momentum $k$ invariant, it is convenient
to consider the generator $U(R_2^{-1})T$, where $R_2$ is the rotation
\eqref{eq:2.6.19}, which also takes $k$ into $\mathcal{P}k$. This commutes with $J_3$,
so
\begin{equation}
U(R_2^{-1})T\ket{k,\sigma}=\zeta_\sigma\ket{k,\sigma}
\tag{2.6.23}\label{eq:2.6.23}
\end{equation}
with $\zeta_\sigma$ another phase. Since $R_2^{-1}\mathcal{T}$ commutes with
the boost \eqref{eq:2.5.45}, and $\mathcal{T}$ commutes with the rotation
$R(\hat{\mathbf{p}})$, operating with $T$ on the state \eqref{eq:2.5.5} gives
\begin{equation}
T\ket{p,\sigma}
=\sqrt{\frac{\kappa}{p^0}}\,
U\left[R(\hat{\mathbf{p}})R_2B\left(\frac{|\mathbf{p}|}{\kappa}\right)\right]
\zeta_\sigma\ket{k,\sigma}.
\tag{2.6.24}\label{eq:2.6.24}
\end{equation}
Using Eq.~\eqref{eq:2.6.21}, this yields finally
\begin{equation}
T\ket{p,\sigma}
=\zeta_\sigma\exp(\pm i\pi\sigma)\ket{\mathcal{P}p,\sigma}.
\tag{2.6.25}\label{eq:2.6.25}
\end{equation}
Again, the top or bottom sign applies according to whether the two-component
of $\mathbf{p}$ is positive or negative, respectively.

\begin{center}
$*\quad *\quad *$
\end{center}

It is interesting that the square $T^2$ of the time-reversal operator has a
very simple action on both massive and massless one-particle states. Using Eq.
\eqref{eq:2.6.18}, and recalling that $T$ is antiunitary, we see that for massive
one-particle states:
\[
\begin{aligned}
T^2\ket{p,\sigma}
&=T\zeta(-1)^{j-\sigma}\ket{\mathcal{P}p,-\sigma}\\
&=\zeta^*(-1)^{j-\sigma}\zeta(-1)^{j+\sigma}\ket{p,\sigma}
\end{aligned}
\]
or in other words
\begin{equation}
T^2\ket{p,\sigma}=(-1)^{2j}\ket{p,\sigma}.
\tag{2.6.26}\label{eq:2.6.26}
\end{equation}
We get the same result for massless particles. If the two-component of
$\mathbf{p}$ is positive then the two-component of $\mathcal{P}\mathbf{p}$ is
negative, and vice-versa, so Eq.~\eqref{eq:2.6.25} gives
\[
\begin{aligned}
T^2\ket{p,\sigma}
&=T\zeta_\sigma\exp(\pm i\pi\sigma)
  \ket{\mathcal{P}p,\sigma}\\
&=\zeta_\sigma^*\exp(\mp i\pi\sigma)
  \zeta_\sigma\exp(\mp i\pi\sigma)\ket{p,\sigma}\\
&=\exp(\mp2\pi i\sigma)\ket{p,\sigma}.
\end{aligned}
\]
As long as $\sigma$ is an integer or half-integer, this can be written
\begin{equation}
T^2\ket{p,\sigma}=(-1)^{2|\sigma|}\ket{p,\sigma}.
\tag{2.6.27}\label{eq:2.6.27}
\end{equation}
By the `spin' of a massless particle, we usually mean the absolute value of
the helicity, so Eq.~\eqref{eq:2.6.27} is the same as Eq.~\eqref{eq:2.6.26}.

This result has an interesting consequence. When $T^2$ acts on any state
$\ket{\Psi}$ of a system of non-interacting particles, either massive or
massless, it yields a factor $(-1)^{2j}$ or $(-1)^{2|\sigma|}$ for each
particle. Hence if the state contains an odd number of particles of
half-integer spin or helicity (plus any number of particles of integer spin or
helicity), we get an overall change of sign
\begin{equation}
T^2\ket{\Psi}=-\ket{\Psi}.
\tag{2.6.28}\label{eq:2.6.28}
\end{equation}
If we now `turn on' various interactions, this result will be preserved,
provided these interactions respect invariance under time-reversal, even if
they do not respect rotational invariance. (For instance, these arguments
will apply even if our system is subjected to arbitrary static gravitational
and electric fields.) Now, suppose that $\ket{\Psi}$ is an eigenstate of the
Hamiltonian. Since $T$ commutes with the Hamiltonian, $T\ket{\Psi}$ will also
be an eigenstate of the Hamiltonian. Is it the same state? If so, then
$T\ket{\Psi}$ can differ from $\ket{\Psi}$ only by a phase
\[
T\ket{\Psi}=\zeta\ket{\Psi},
\]
but then
\[
T^2\ket{\Psi}=T(\zeta\ket{\Psi})
=\zeta^*T\ket{\Psi}=|\zeta|^2\ket{\Psi}=\ket{\Psi},
\]
in contradiction with Eq.~\eqref{eq:2.6.28}. We see that any energy eigenstate
$\ket{\Psi}$ satisfying Eq.~\eqref{eq:2.6.28} must be degenerate, with another
eigenstate of the same energy. This is known as a `Kramers
degeneracy'~\hyperref[ref:2.10]{[10]}. Of course, this conclusion is trivial if the
system is in a rotationally invariant environment, because the total
angular-momentum $j$ of any state of this system would have to be a
half-integer, and there would therefore be $2j+1=2,4,\ldots$ degenerate states.
The surprising result is that at least a two-fold degeneracy persists even if
rotational invariance is perturbed by external fields, such as electrostatic
fields, as long as these fields are invariant under $T$. In particular, if
any particle had an electric or gravitational dipole moment then the
degeneracy among its $2j+1$ spin states would be entirely removed in a static
electric or gravitational field, so such dipole moments are forbidden by
time-reversal invariance.

For the sake of completeness, it should be mentioned that $P$ and $T$ can
have more complicated effects on multiplets of particles with the same mass.
This possibility will be considered in Appendix C of this chapter. No
physically relevant examples are known.
